\documentclass[11pt]{article}
\usepackage[margin=1in]{geometry}
\usepackage{amsmath,amssymb,amsthm}
\usepackage{hyperref}
\usepackage[numbers,sort&compress]{natbib}

\title{The Algebraic Core of Constraint Grammar as Set Theoretic Estimation}
\author{Project \texttt{thejeb}}
\date{\today}

\newtheorem{claim}{Claim}
\newtheorem{definition}{Definition}

\begin{document}
\maketitle

\begin{abstract}
This note does not argue for Constraint Grammar (CG) as a parsing technology.
Instead it isolates an algebraic question suggested by CG: when a grammar starts
with a universe of candidate analyses and rules eliminate candidates, is that
algebra the same as set theoretic estimation (STE)?  For the hard, purely
eliminative fragment the answer is yes.  A CG rule is exactly an STE property
set, the surviving analyses are exactly the STE feasibility set, adding rules is
exactly antitone refinement, fairness is preservation of the intended parse, and
unambiguous successful analysis is ideal information.  The Lean module
\texttt{Ste.ConstraintGrammar} verifies this identification by giving an
explicit equivalence between hard-constraint grammars and STE property-set
families.  The equivalence is deliberately scoped: CG operations that genuinely
change the candidate universe are not isomorphic to static STE unless one first
embeds them into a fixed enriched universe or treats them as staged maps between
STE problems.
\end{abstract}

\section{Question}
Fix a sentence $s$ only to obtain a finite universe $\mathcal P_s$ of possible
complete analyses.  The linguistic details of $\mathcal P_s$ are not the point.
The algebraic question is:
\[
  \text{candidate universe} + \text{hard eliminating constraints}
  \quad\cong?\quad
  \text{STE solution space} + \text{property sets}.
\]
Combettes's STE formulation represents each item of information as a property
set $S_i\subseteq\Xi$ and estimates by the intersection $\bigcap_i S_i$
\citep{combettes1993}.  Classical CG is presented as a rule-based framework for
reducing locally possible readings in context \citep{karlsson1990constraint,
karlsson1995constraint}; practical CG systems include richer operations and
rule languages \citep{bick2015cg3}.  The proposed connection concerns the
common hard-constraint core, not every operational feature of every parser.


\section{Abstract STE algebra}
The more general algebraic description is independent of parsing.  Let $B$ be a
space of problem instances, $X$ a space of candidate solutions, and $C$ a space
of criteria.  A crisp STE algebra is an interpretation map
\[
  \alpha : B \to C \to \mathcal P(X).
\]
For each problem instance $b\in B$, the criterion $c\in C$ denotes the property
set $\alpha(b,c)$ of solutions satisfying that criterion.  The STE solution
operator is the meet in the powerset lattice:
\[
  \operatorname{Sol}_\alpha(b) = \bigwedge_{c\in C} \alpha(b,c)
  = \bigcap_{c\in C} \alpha(b,c).
\]
Thus an instance is STE-solvable when this meet is nonempty, and it is
identified by $x\in X$ when the meet is the singleton $\{x\}$.  Soundness or
fairness for a proposed true solution $x$ means $x\in\alpha(b,c)$ for every
criterion $c$, which immediately implies solvability.

The Lean module \texttt{Ste.Algebra} packages this as
\texttt{STE.Algebra Problem Solution Criterion}.  It records the same structure
as the curried map \texttt{Problem -> Criterion -> Set Solution}, proves the
membership criterion for the meet, defines \texttt{Solvable}, \texttt{SoundFor},
and \texttt{Identifies}, and reuses the STE theorems showing that soundness
implies solvability and singleton feasibility implies uniqueness.  In this
language, the CG result is simply the special case where the problem instance is
a fixed sentence-analysis task and the criteria are hard rules.

\section{The isomorphism for hard constraints}
Let $P$ be a type or set of candidate parses and let $R$ be a type or set of
rules.  A hard CG interpretation is a function
\[
  A : R \to \mathcal P(P),
\]
where $A(r)$ is the set of parses accepted, or surviving, rule $r$.  The set of
surviving parses is
\[
  \operatorname{Parses}(A) = \bigcap_{r\in R} A(r).
\]
An STE presentation over the same universe and index set is also exactly a
function
\[
  S : R \to \mathcal P(P),
\]
with feasibility set
\[
  \operatorname{Feas}(S) = \bigcap_{r\in R} S(r).
\]
Thus the data are not merely analogous: after abstracting rules by their
extension, they are definitionally the same data.  In Lean this is recorded as
an equivalence
\[
  \texttt{Grammar Parse Rule} \simeq (\texttt{Rule} \to \texttt{Set Parse}),
\]
implemented by \texttt{STE.ConstraintGrammar.grammarEquivSTE}.  The theorem
\texttt{parses\_eq\_feasibilitySet} then states that the parse survivor set is
exactly the STE feasibility set under this equivalence.

\begin{claim}[Hard CG $\cong$ crisp STE]
For fixed candidate universe $P$ and rule index set $R$, extensional hard CGs
are isomorphic to crisp STE property-set families indexed by $R$.
\end{claim}

\section{Consequences transported from STE}
Because the algebraic data are the same, the basic STE theorems transfer
without new linguistic assumptions.

\paragraph{Membership.}
A parse $p$ survives the grammar iff it survives every rule:
\[
  p\in\operatorname{Parses}(A) \iff \forall r\in R,\ p\in A(r).
\]

\paragraph{Monotonicity.}
If $R_0\subseteq R_1$, then applying all rules in $R_1$ can only shrink the
survivor set relative to applying $R_0$:
\[
  \bigcap_{r\in R_1} A(r) \subseteq \bigcap_{r\in R_0} A(r).
\]
This is the algebraic invariant behind rule accumulation.  It does not say that
rules are linguistically correct; it says only that hard constraints refine by
intersection.

\paragraph{Fairness.}
For an intended parse $p_\star$, a rule is fair when $p_\star\in A(r)$.  A fair
grammar preserves $p_\star$ and therefore has a nonempty survivor set.  This is
the CG version of STE consistency from fair information.

\paragraph{Idealness.}
A grammar is ideal for $p_\star$ when
\[
  \operatorname{Parses}(A)=\{p_\star\}.
\]
Then any surviving parse is the intended parse.  Ambiguity is simply nonempty,
non-singleton feasibility; failure is empty feasibility.


\section{Reduction, priority, and contradiction}
Once STE is viewed as meet in a powerset lattice, every new property set has an
exact set-valued reduction.  If $F_C$ is the current partial feasible set and a
new criterion $c$ contributes property set $S_c$, then
\[
  F_{C\cup\{c\}} = F_C \cap S_c,
  \qquad
  R_c(F_C) = F_C \setminus (F_C\cap S_c).
\]
The reduction is nonnegative in the order-theoretic sense: the new feasible set
is a subset of the old one, and the removed candidates are a subset of the old
feasible set.  The engineer's question ``which criterion shrinks fastest?'' is
therefore a question about choosing or estimating a score
\[
  r(b,c) \approx \operatorname{size}\bigl(R_c(F_C)\bigr).
\]
Depending on the application, $r$ might be a finite cardinality decrease, a
measure decrease, an expected posterior reduction, an entropy estimate, or an
AEP/typical-set estimate.  Carlson's AEP move can be read in exactly this way:
use probabilistic typicality to identify property sets that should eliminate a
large part of the key/message universe while preserving the true instance.

The same accounting also distinguishes productive reduction from malformation.
A strong criterion may shrink rapidly while remaining fair.  A malformed pair
of criteria may be disjoint:
\[
  S_c \cap S_d = \varnothing.
\]
If both are enforced, the global meet is empty.  This is not successful
identification; it is inconsistency.  In the algebra, emptiness means at least
one criterion, candidate generator, or modeling assumption is wrong relative to
the intended solution.  The Lean module \texttt{Ste.Reduction} reifies this by
introducing exact \texttt{reductionSet}s, proving that adding a criterion
shrinks the current partial feasible set, and proving that incompatible
criteria imply non-solvability.

\section{Where the isomorphism stops}
The equivalence is intentionally limited to extensional hard constraints on a
fixed universe.  Some CG implementations include operations such as add, map,
replace, subreadings, traces, dynamic cohorts, or staged rule side effects
\citep{bick2015cg3}.  Such operations are not literally static intersections if
they change the universe on which later rules operate.  There are two ways to
recover an STE account:
\begin{enumerate}
\item \textbf{Enriched universe:} define $P$ broadly enough to contain complete
  before-and-after analyses, mapped tags, and latent features.  Each operation
  is then represented extensionally as a predicate on complete histories or
  complete enriched parses.
\item \textbf{Staged STE:} model each operation as a map from one STE problem to
  another, rather than as a single property set in one fixed feasibility
  problem.
\end{enumerate}
Therefore the precise theorem is not ``all Constraint Grammar implementations
are STE.''  The verified theorem is narrower and stronger: the algebra of fixed
candidate-universe hard constraint elimination is isomorphic to crisp STE.

\section{Topological reading}
For a fixed finite sentence analysis space, give $P$ the discrete topology.
Every hard rule property set is closed.  If $P$ is finite, it is compact in the
usual discrete setting, so the existing topological STE theorem applies:
finite consistency of closed rule constraints gives nonempty global feasibility.
This is mathematically modest for a finite, explicitly enumerated rule set, but
it is useful as a conceptual bridge from Carlson's finite/discrete topological
STE perspective \citep{carlson2012} to non-cryptographic symbolic analysis.

\section{Lean status}
The mechanized contribution is \texttt{lean/Ste/ConstraintGrammar.lean}.  It
contains:
\begin{itemize}
\item \texttt{Grammar}, a structure with field \texttt{accepts : Rule -> Set Parse};
\item \texttt{parses}, defined as \texttt{STE.feasibilitySet accepts};
\item \texttt{grammarEquivSTE}, an explicit equivalence between grammars and
  STE property-set families;
\item transport lemmas for membership, fairness/nonemptiness, ideal
  identification, partial rule sets, and antitone refinement.
\end{itemize}
This verifies the algebraic isomorphism in the crisp hard-constraint setting
and records the boundary where non-eliminative dynamics require additional
modeling.

\bibliographystyle{plainnat}
\bibliography{../../refs}

\end{document}
